\section{Introdução: A Encruzilhada Ética da Simulação}

A ascensão dos gêmeos digitais na odontologia contemporânea deflagra um labirinto de desafios éticos, legais e regulatórios que exigem resolução sistemática, sob pena de a inovação tecnológica colidir frontalmente com os direitos fundamentais do paciente. Ao contrário das inovações analógicas precedentes — focadas predominantemente em materiais restauradores ou cinemática de instrumentação —, os gêmeos digitais operam na interseção crítica entre biometria, predição algorítmica e intervenção clínica. Eles aglutinam, em uma única entidade virtual, dados pessoais sensíveis, simulações biomecânicas de alto risco e o monitoramento longitudinal do indivíduo.

A literatura acadêmica de ponta tem soado alarmes sobre a insuficiência das matrizes jurídicas vigentes. O tratado de Springer \cite{springer2026realising} sobre as questões Éticas, Legais e Sociais (ELSI) dos gêmeos digitais na saúde evidencia que os arcabouços atuais, desenhados para proteger Prontuários Eletrônicos do Paciente (PEP) estáticos e equipamentos médicos convencionais, são obsoletos perante sistemas autônomos de \textit{Continuous Learning}. No Brasil, um marco fundacional foi estabelecido em fevereiro de 2026 com a promulgação da \textbf{ABNT NBR ISO/IEC 3173} \cite{abinc2026brasil}. Embora não legisle sobre bioética, a norma proveu o ancoradouro terminológico necessário para que os tribunais e comitês de ética comecem a debater as fronteiras da tecnologia com rigor conceitual.

Este capítulo disseca esta encruzilhada. As seções \ref{sec:privacidade} e \ref{sec:regulacao} examinam a tessitura da Lei Geral de Proteção de Dados (LGPD) e o arcabouço da ANVISA sob o prisma dos gêmeos digitais. As instâncias subsequentes escrutinam o dilema da responsabilidade civil (\ref{sec:responsabilidade}), as disparidades no acesso (\ref{sec:equidade}), os vieses algorítmicos intrínsecos à simulação (\ref{sec:etica-simulacao}) e os conflitos sobre a titularidade dos dados (\ref{sec:propriedade-intelectual}), culminando em uma proposta de governança bioética (\ref{sec:conclusao}).
